\section{Introduction}
\label{sec:intro}
%% Packet W-11 (first half; written last, from the frozen results sections, per
%% design/paper_plan.md and paper/narrative.md). Readability pass 2026-09-22: rewritten for a
%% glass-physics reader; claims unchanged.

\subsection{The problem}

Anderson called the nature of glass and the glass transition probably the deepest and most
interesting unsolved problem in solid-state theory \cite{anderson-1995}. One face of that problem is sixty years old and
unusually concrete: no universally accepted static order parameter for glass has emerged. A glass and a liquid
can agree in density, energy, and every routinely monitored one-time observable, yet behave
differently \cite{debenedetti-stillinger-2001}; conventional low-order static correlations show no
clear signature of the transition, and the candidate structural measures that do grow remain
contested \cite{berthier-biroli-2011,cavagna-2009,tarjus-2011,royall-williams-2015}. A second face has sat in the laboratory record since 1963: Kovacs's
crossover experiment \cite{kovacs-1963}. Equilibrate a polymer glass hot, quench it, age it, then
jump the temperature up and watch its volume relax. At some moment the volume \emph{equals} its
new equilibrium value, and does not stay there: it overshoots through a non-monotone hump before
returning. At the crossing moment the sample is indistinguishable from equilibrium in the
monitored observable and distinguishable from it in its future, a point made explicit at the
level of configurations in simulation \cite{mossa-sciortino-2004}. The standard partial fix, the
Tool--Narayanaswamy--Moynihan fictive temperature \cite{tool-1946,narayanaswamy-1971,moynihan-1976},
adds one memory variable; richer protocols expose the limits of any single such variable and
have motivated multiparameter models \cite{ritland-1956,kahr-1979,hodge-1994,mckenna-simon-2017}. Memory effects
of this kind, in which a material's response depends on its history in a way no monitored
one-time observable records, are now studied across many classes of matter \cite{keim-2019}.

\subsection{The question, made precise}

Both faces are the same question. Take a declared readout of the system (here, the energy of a
spin chain; in the laboratory, the volume) and ask: \emph{is that readout a closed description of
the dynamics?} A readout is closed over a horizon $\tau$ if its present value determines the
statistics of its value $\tau$ steps later, so that a law can be written in the readout's
vocabulary without loss, the informational closure of \cite{bertschinger-2006,pfante-2014}. In Markov-chain language the readout is closed at horizon $\tau$ exactly
when the $\tau$-step readout transition probabilities agree across every populated microstate with
the same readout, which for a fully supported distribution is lumpability of the readout
\cite{burke-rosenblatt-1958,kemeny-snell-1960}. In this paper a static order-parameter candidate is taken in one operational sense only: a
function of the declared readout that separates two preparations whose futures differ. Any
function of the readout that fails to separate them (a constant does so trivially) cannot carry
the information the futures carry, and that is the sense in which GT3 below forecloses the
search: not that no coarsening of the energy is closed, but that none of them tells the exhibited
pair apart. That is narrower than the field's usage. The Kovacs hump is a direct demonstration
that the monitored volume is not closed.

Once the question is put this way, three further questions become computable on a finite model.
How far from closed is the readout, in nats? Can one exhibit a concrete certificate of the
failure, that is, two preparation histories that agree on the readout now and differ in their
futures? And is there \emph{any} function of the readout that recovers the missing predictive
distinction, or must it be carried by a genuinely new coordinate, and at what cost?

\subsection{What this paper does}

We answer those questions on a class of models small enough that the certificates are exact
ratios of integers wherever the artifact declares them so, and deterministic float64 elsewhere:
kinetically constrained models
\cite{ritort-sollich-2003,garrahan-sollich-toninelli-2011}, with the East model
\cite{jackle-eisinger-1991} as the substrate on which everything is built and calibrated, and
the Fredrickson--Andersen model \cite{fredrickson-andersen-1984} and a Bouchaud-type trap family
\cite{bouchaud-1992,monthus-bouchaud-1996} as out-of-sample transfer targets. At the system sizes
used ($N \le 10$ spins for certificates), transition kernels, stationary measures, witness pairs,
and zero controls are ratios of integers; the closure-deficit grid of GT1 is deterministic
float64 and is labeled as such in its artifact, with its margin far from the rounding scale. The
price is that every result is a statement about a declared finite class of models, sizes,
observables, and protocols, and we state that scope in every claim.

The bookkeeping that keeps those scopes honest is borrowed from a certificate-based closure
calculus, Six Birds Theory \sbtcite{sbt-primer,sbt-foundations-i,sbt-foundations-ii,%
sbt-foundations-iii,sbt-foundations-iv,sbt-paper-a,sbt-paper-b,sbt-paper-c}. The reader does not
need that corpus to follow this paper: \S\ref{sec:methods} defines every object we use in
ordinary probabilistic terms, and the framework's own labels appear only where a claim is filed in
its vocabulary, always with a plain gloss. Readers who want the formal anchors will find each
claim's corpus reference in the claims ledger shipped with the paper.

The results, each at its declared scope (\S\ref{sec:gt1}--\S\ref{sec:gt7}), read as follows.

\begin{itemize}
  \item \textbf{GT1, the energy readout is not closed.} For an aging East chain, the closure
    deficit of the energy readout (the conditional mutual information between the microstate now
    and the energy later, given the energy now) exceeds the same quantity at constrained
    equilibrium by a certified factor of at least $3/2$. A correction to our own design
    expectations came with it: the equilibrium value is provably \emph{not} zero, because
    closure depends on where the mobile excitations sit, not on whether the chain is stationary.
  \item \textbf{GT2, the Kovacs witness.} The Kovacs protocol on East yields an exact pair of
    preparation histories with the same present mean energy and different future mean energies
    under the declared continuation.
  \item \textbf{GT3, the order-parameter search is foreclosed on the declared class.} Every one of
    the $512$ yes/no functions definable from the declared energy readout agrees on a witness
    pair that a finer readout of the same microstate (the occupation of the site next to the
    wall) separates. On that class, no static order parameter built from the energy readout can
    do the job; whatever does the job must carry information the energy readout does not.
  \item \textbf{GT4, the kill-test panel.} The classifier that accepts the Kovacs memory
    candidate correctly rejects five constructed look-alike regimes, schematic fixtures for a
    protocol artifact, dissipation, an explicitly latent variable, an equilibrium control, and an
    incomplete continuation catalog.
  \item \textbf{GT5, the cost of repair.} One added scalar coordinate resolves the witness pair,
    and a fictive-temperature-like index (the index of the grid point whose equilibrium mean energy is closest to each
    microstate's energy, averaged over the history) is one of three that do. The
    certificate shows the repair mechanism on this pair; it does not prove that one scalar is
    insufficient in general.
  \item \textbf{GT6, the memory layer is strictly richer.} The description of histories by
    future-equivalence strictly extends the declared energy-based description, certified for the
    audited shell (the frozen set of eligibility conditions, \S\ref{sec:methods}) by an explicit
    gate table and ablation panel.
  \item \textbf{GT7, aging constants transfer only partly.} Six aging readouts were frozen on
    East, predictions for FA and the trap family were pre-registered, and the targets were then
    run and scored mechanically: 8 of 22 scored predictions passed, 14 failed on the same side,
    and the trap family produced no Kovacs hump at all, for a reason we prove.
\end{itemize}

Two further results come from the formalization effort (\S\ref{sec:lean}). A vendored core of
the framework's theorems is machine-checked in the Lean proof assistant, and the attempt to encode
the numeric Kovacs witness in that core produced a machine-checked impossibility theorem instead,
which locates exactly which certificates can live in that core and why.

\subsection{How to read the claims}

Four contributions, in the order we would defend them. First, the \textbf{discipline}: every
claim quantifies over a declared finite class, is pinned by content hash to the artifact that
supports it, is accompanied by controls, and is filed together with what it does \emph{not}
assert; honest negatives are results, and the discipline forced several of our own claims to
shrink. Second, the \textbf{result complex} itself: seven claims, four bridges, and six honest
negatives, each hash-pinned to its artifacts. The bracketed labels throughout, such as [GT1], are
machine-checked citations into the claims ledger, not decoration. Third, the
\textbf{freeze-and-transfer protocol with its actual scoreboard}: the failures are results, with
a clean one-sided pattern and a proven structural explanation for the empty column. Fourth, the
\textbf{Lean-checked core and the scope-boundary finding}.

What this paper is not: it proves nothing about all glasses, the continuum, or the ideal-glass
transition; it makes no irreversibility claims; and its field-facing sentence, ``no static order
parameter,'' is licensed only over the declared class it is proved on. The full list of nonclaims
is printed, unabridged, in \S\ref{sec:discussion}.
